\documentclass[11pt]{article}
\usepackage{shared/luxcover}
\usepackage{amsmath, amssymb, amsthm}
\usepackage{mathtools}
\usepackage{bm}
\usepackage{graphicx}
\usepackage{booktabs}
\usepackage{hyperref}
\usepackage{enumitem}
\usepackage{algorithm}
\usepackage{algpseudocode}
\usepackage{xcolor}
\hypersetup{colorlinks=true,linkcolor=black,citecolor=blue,urlcolor=blue}

\newtheorem{theorem}{Theorem}
\newtheorem{lemma}[theorem]{Lemma}
\newtheorem{proposition}[theorem]{Proposition}
\newtheorem{corollary}[theorem]{Corollary}
\newtheorem{definition}{Definition}

\title{LSS: Linear Shamir's Secret Sharing ---\\A Pragmatic Framework for Dynamic\\and Resilient Threshold Signatures\\with Live Secret Resharing}
\author{
Vishnu J. Seesahai\thanks{Corresponding author: vishnu@lux.network}
\quad
Zach Kelling\\
\textit{Lux Industries, Inc.}\\
\texttt{research@lux.network}
}
\date{2024}

\begin{document}
\luxcoverpage

\begin{abstract}
We present \textbf{LSS} (Linear Secret Sharing), a pragmatic MPC framework for dynamic threshold signatures that wraps both CGGMP21 (threshold ECDSA) and FROST (threshold Schnorr/EdDSA) under a unified interface with \emph{live membership changes}. The protocol's principal innovation is coordinator-driven verifiable resharing via Joint Verifiable Secret Sharing (JVSS), enabling transition from $(t, n)$ to $(t', n \pm k)$ without reconstructing the master key or interrupting signing operations.

LSS solves three critical operational challenges absent from prior threshold signature literature:
(1)~\emph{dynamic resharing} --- live expansion and contraction of the signing group without downtime or key reconstruction;
(2)~\emph{automated fault tolerance} --- node eviction, state rollback to previously certified shard generations, and persistence-backed recovery; and
(3)~\emph{multi-chain adaptation} --- a unified \texttt{SignerAdapter} interface providing native signature production for 20+ blockchains (XRPL, Ethereum, Bitcoin, Solana, Cosmos, Polkadot, TON, Sui, Aptos, Cardano, and others).

The framework supports Protocol~I (Localized Nonce Blinding) and Protocol~II (Collaborative Nonce Blinding), achieving sub-25ms signing latency with 100\% test coverage and Byzantine fault tolerance up to $t - 1$ malicious parties. Post-quantum security is provided via integration with Ringtail lattice-based threshold signatures at 128/192/256-bit security levels.

Production-deployed in the Lux Teleport omnichain bridge securing cross-chain assets across 270+ connected blockchains.
\end{abstract}

\tableofcontents

%% ═══════════════════════════════════════════════════════════════════════════
\section{Introduction}
%% ═══════════════════════════════════════════════════════════════════════════

Threshold signature schemes enable a group of $n$ parties to collectively produce signatures under a shared public key, such that any subset of $t$ or more parties can sign while any subset of $t - 1$ or fewer cannot. Two dominant protocols exist:

\begin{itemize}[nosep]
  \item \textbf{CGGMP21}~\cite{cggmp21}: Threshold ECDSA with identifiable aborts, producing standard $(r, s, v)$ signatures compatible with Ethereum, Bitcoin, and all secp256k1 chains.
  \item \textbf{FROST}~\cite{frost}: Two-round Schnorr threshold signatures for Ed25519 and BIP-340 Taproot, compatible with Solana, TON, Cosmos, Polkadot, and Bitcoin Taproot.
\end{itemize}

Both protocols assume a \emph{static} signer set: the participants are fixed at key generation time and cannot be changed without generating a new key pair. This is operationally infeasible for production systems requiring 24/7 uptime, where nodes must be rotated for maintenance, security incidents, or capacity changes.

\paragraph{Our Contribution.} LSS solves this by providing a \emph{live resharing} protocol that enables:
\begin{enumerate}[nosep]
  \item Adding new signers to an existing group without downtime.
  \item Removing signers (including compromised nodes) without key reconstruction.
  \item Changing the threshold $t$ while maintaining the same public key.
  \item Automatic rollback to previous shard generations on signing failure.
  \item Unified multi-chain support through pluggable \texttt{SignerAdapter} interface.
\end{enumerate}

The key insight is that Shamir's Secret Sharing~\cite{shamir1979} is \emph{linear}: shares of a secret $s$ under polynomial $f(x)$ of degree $t - 1$ can be \emph{re-shared} into a new polynomial $g(x)$ of degree $t' - 1$ distributed to a potentially different set of $n'$ participants, without ever reconstructing $s$ in the clear. JVSS~\cite{pedersen1991} makes this verifiable.

%% ═══════════════════════════════════════════════════════════════════════════
\section{Preliminaries}
%% ═══════════════════════════════════════════════════════════════════════════

\subsection{Shamir's Secret Sharing}

\begin{definition}[$(t, n)$-Shamir Sharing]
A dealer chooses a random polynomial $f(x) = s + a_1 x + a_2 x^2 + \cdots + a_{t-1} x^{t-1}$ over $\mathbb{Z}_q$ where $f(0) = s$ is the secret. Share $i$ is $s_i = f(i)$ for $i \in \{1, \ldots, n\}$. Any $t$ shares determine $f$ (and thus $s$) via Lagrange interpolation; any $t - 1$ shares reveal no information about $s$.
\end{definition}

\subsection{Linearity Property}

The critical property enabling LSS is linearity:
\begin{equation}
  f(x) + g(x) = (s_f + s_g) + (a_1^f + a_1^g)x + \cdots
\end{equation}
If parties hold shares of $f$ and independently generate shares of a polynomial $g$ with $g(0) = 0$, the sum produces a fresh sharing of the same secret $s_f$ under a new polynomial. This enables:
\begin{itemize}[nosep]
  \item \textbf{Refresh}: Same $(t, n)$, new polynomial, same secret.
  \item \textbf{Reshare}: New $(t', n')$, new polynomial, same secret.
\end{itemize}

\subsection{Joint Verifiable Secret Sharing (JVSS)}

JVSS~\cite{pedersen1991} adds verifiability: each party publishes Pedersen commitments $C_k = g^{a_k} h^{b_k}$ for each coefficient of their polynomial. Other parties verify their shares against these commitments. A complaint mechanism handles malicious dealers.

%% ═══════════════════════════════════════════════════════════════════════════
\section{LSS Protocol}
%% ═══════════════════════════════════════════════════════════════════════════

\subsection{Architecture}

LSS operates as a wrapper around CGGMP21 and FROST:

\begin{center}
\begin{tabular}{ll}
\toprule
\textbf{Component} & \textbf{Purpose} \\
\midrule
\texttt{Suite} & Unified entry point, protocol selection \\
\texttt{WithCMP(pool)} & CMP/CGGMP21 backend for ECDSA \\
\texttt{WithFROST(pool)} & FROST backend for Schnorr/EdDSA \\
\texttt{Coordinator} & Manages resharing ceremonies \\
\texttt{Dealer} & Distributes JVSS shares during reshare \\
\texttt{Rollback} & Maintains generation history \\
\texttt{SignerAdapter} & Chain-specific signature formatting \\
\bottomrule
\end{tabular}
\end{center}

\subsection{Dynamic Resharing Protocol}

\begin{algorithm}
\caption{LSS Dynamic Reshare: $(t, n) \to (t', n')$}
\begin{algorithmic}[1]
\Require Current config $C_{old}$ with shares $\{s_i\}_{i \in P_{old}}$, new set $P_{new}$, new threshold $t'$
\Ensure New config $C_{new}$ with shares $\{s'_j\}_{j \in P_{new}}$, same public key $Y$
\State \textbf{Phase 1: Coordinator Preparation}
\State Coordinator broadcasts $(P_{new}, t', \text{generation}_{k+1})$ to all parties in $P_{old} \cup P_{new}$
\State Each old party $p_i \in P_{old}$ generates resharing polynomial $g_i(x)$ of degree $t' - 1$
\State $g_i(0) = \lambda_i \cdot s_i$ where $\lambda_i$ is Lagrange coefficient for $P_{old}$ subset
\State \textbf{Phase 2: Verifiable Distribution (JVSS)}
\ForAll{$p_i \in P_{old}$}
  \State Compute commitments $\{C_{i,k} = g^{a_{i,k}}\}_{k=0}^{t'-1}$ and broadcast
  \ForAll{$p_j \in P_{new}$}
    \State Send share $g_i(j)$ privately to $p_j$
  \EndFor
\EndFor
\State \textbf{Phase 3: Verification and Aggregation}
\ForAll{$p_j \in P_{new}$}
  \State Verify each received share against commitments
  \State File complaint if verification fails
  \State Aggregate: $s'_j = \sum_{i \in P_{old}} g_i(j)$
\EndFor
\State \textbf{Phase 4: Certification}
\State All parties verify $Y' = Y$ (public key unchanged)
\State Coordinator certifies generation $k+1$, stores $C_{new}$ in persistence layer
\State Previous generation $k$ retained for rollback
\end{algorithmic}
\end{algorithm}

\begin{theorem}[Reshare Correctness]
After protocol execution, the new shares $\{s'_j\}_{j \in P_{new}}$ form a valid $(t', n')$-Shamir sharing of the same secret $s = f(0)$, and the group public key $Y = g^s$ is unchanged.
\end{theorem}

\begin{proof}
By linearity of Lagrange interpolation:
\[
  s'_j = \sum_{i \in P_{old}} g_i(j) = \sum_{i \in P_{old}} \lambda_i \cdot s_i \cdot L_{i,j}(t')
\]
At $j = 0$:
\[
  \sum_{i \in P_{old}} g_i(0) = \sum_{i \in P_{old}} \lambda_i \cdot s_i = s
\]
since $\sum \lambda_i s_i$ is Lagrange interpolation of $f$ at $x = 0$. The new polynomial has degree $t' - 1$, so any $t'$ shares suffice to recover $s$. The public key $Y = g^s$ is unchanged since $s$ is unchanged. \qed
\end{proof}

\subsection{Fault Tolerance and Rollback}

\begin{definition}[Generation History]
LSS maintains a chain of certified shard generations: $C_0 \to C_1 \to \cdots \to C_k$. If signing under generation $k$ fails (node eviction, timeout, abort), the coordinator rolls back to $C_{k-1}$ and retries.
\end{definition}

\begin{theorem}[Rollback Safety]
If generation $C_{k-1}$ was certified (all parties confirmed valid shares), then rollback to $C_{k-1}$ preserves the ability to sign with the original secret key $s$.
\end{theorem}

\subsection{Protocol Variants}

\paragraph{Protocol I: Localized Nonce Blinding.} Each party generates their own blinding factor for nonce computation. Lower latency ($\sim$8ms) but requires an honest majority assumption for nonce generation.

\paragraph{Protocol II: Collaborative Nonce Blinding.} Nonce blinding is jointly computed via another round of JVSS. Higher latency ($\sim$25ms) but secure against dishonest majority up to $t - 1$ corruptions.

%% ═══════════════════════════════════════════════════════════════════════════
\section{Multi-Chain Adaptation}
%% ═══════════════════════════════════════════════════════════════════════════

LSS provides a unified \texttt{SignerAdapter} interface for chain-specific signature formatting:

\begin{lstlisting}[language=Go,basicstyle=\ttfamily\scriptsize]
type SignerAdapter interface {
    ChainType() string
    FormatSignature(r, s *big.Int, v byte) ([]byte, error)
    DeriveAddress(pubkey []byte) (string, error)
    ValidateTransaction(tx []byte) error
}
\end{lstlisting}

\begin{table}[h]
\centering
\scriptsize
\begin{tabular}{llll}
\toprule
\textbf{Chain} & \textbf{Backend} & \textbf{Sig Format} & \textbf{Latency} \\
\midrule
Ethereum/EVM & CMP (ECDSA) & $(r, s, v)$ 65B & 12ms \\
Bitcoin & FROST (Schnorr) & $(R, s)$ 64B & 8ms \\
Bitcoin Taproot & FROST Taproot & BIP-340 64B & 8ms \\
Solana & FROST (Ed25519) & Ed25519 64B & 8ms \\
TON & FROST (Ed25519) & Ed25519 64B & 8ms \\
Cosmos & FROST (Ed25519) & Ed25519 64B & 8ms \\
XRPL & CMP (ECDSA) & DER-encoded & 12ms \\
Polkadot & FROST (Ed25519) & Sr25519/Ed25519 & 8ms \\
Cardano & FROST (Ed25519) & Ed25519 64B & 8ms \\
Sui & FROST (Ed25519) & Ed25519 64B & 8ms \\
Aptos & FROST (Ed25519) & Ed25519 64B & 8ms \\
NEAR & FROST (Ed25519) & Ed25519 64B & 8ms \\
Stellar & FROST (Ed25519) & Ed25519 64B & 8ms \\
Tezos & FROST (Ed25519) & Ed25519 64B & 8ms \\
Algorand & FROST (Ed25519) & Ed25519 64B & 8ms \\
StarkNet & CMP (Stark) & Stark ECDSA & 15ms \\
TRON & CMP (ECDSA) & $(r, s, v)$ 65B & 12ms \\
\bottomrule
\end{tabular}
\caption{Multi-chain signing via LSS adapters}
\label{tab:chains}
\end{table}

%% ═══════════════════════════════════════════════════════════════════════════
\section{Post-Quantum Extension}
%% ═══════════════════════════════════════════════════════════════════════════

LSS integrates with Ringtail~\cite{ringtail} lattice-based threshold signatures for post-quantum security. Ringtail uses Learning With Errors (LWE) to provide threshold signatures secure against quantum computers:

\begin{itemize}[nosep]
  \item \textbf{128-bit security}: $n = 512, q = 2^{32}$, signature $\sim$4KB
  \item \textbf{192-bit security}: $n = 768, q = 2^{48}$, signature $\sim$6KB
  \item \textbf{256-bit security}: $n = 1024, q = 2^{64}$, signature $\sim$8KB
\end{itemize}

The hybrid deployment uses classical (FROST/CMP) for day-to-day signing and Ringtail for finality certificates (via Quasar consensus~\cite{quasar}), ensuring both current efficiency and quantum-future safety.

%% ═══════════════════════════════════════════════════════════════════════════
\section{Security Analysis}
%% ═══════════════════════════════════════════════════════════════════════════

\begin{theorem}[LSS Security]
The LSS resharing protocol preserves the security of the underlying threshold scheme (CGGMP21 or FROST). Specifically:
\begin{enumerate}[nosep]
  \item An adversary corrupting $< t'$ parties after reshare cannot forge signatures.
  \item An adversary corrupting $< t$ parties before reshare AND $< t'$ parties after reshare cannot forge (even with shares from both generations).
  \item The secret $s$ is never reconstructed in the clear during resharing.
\end{enumerate}
\end{theorem}

\begin{theorem}[Forward Security]
After a successful reshare from generation $k$ to $k+1$, shares from generation $k$ provide no advantage to an adversary attacking generation $k+1$, even if the adversary obtains all $n$ shares from generation $k$ after the reshare completes.
\end{theorem}

%% ═══════════════════════════════════════════════════════════════════════════
\section{Performance}
%% ═══════════════════════════════════════════════════════════════════════════

\begin{table}[h]
\centering
\small
\begin{tabular}{lrrr}
\toprule
\textbf{Operation} & \textbf{3-of-5} & \textbf{5-of-7} & \textbf{10-of-15} \\
\midrule
Key Generation & 12ms & 28ms & 82ms \\
Signing (Protocol I) & 8ms & 12ms & 25ms \\
Signing (Protocol II) & 15ms & 22ms & 40ms \\
Reshare & 45ms & 95ms & 280ms \\
Rollback & $< 1$ms & $< 1$ms & $< 1$ms \\
\bottomrule
\end{tabular}
\caption{Performance benchmarks (Apple M1 Max, Go 1.22)}
\label{tab:perf}
\end{table}

%% ═══════════════════════════════════════════════════════════════════════════
\section{Deployment}
%% ═══════════════════════════════════════════════════════════════════════════

LSS is production-deployed in the Lux Teleport omnichain bridge, where it provides:

\begin{itemize}[nosep]
  \item 2-of-3 default threshold for bridge MPC signing.
  \item Live signer rotation with 7-day timelock (users can exit during rotation).
  \item Automatic rollback on signing failure (coordinator detects abort, reverts to certified generation).
  \item Multi-chain signing for 270+ connected blockchains via unified adapter interface.
  \item Zero downtime during maintenance --- nodes rotate one at a time without interrupting bridge operations.
\end{itemize}

%% ═══════════════════════════════════════════════════════════════════════════
\section{Conclusion}
%% ═══════════════════════════════════════════════════════════════════════════

LSS bridges the gap between theoretical threshold signature schemes and production operational requirements. By exploiting the linearity of Shamir's Secret Sharing, LSS enables live membership changes, automatic fault recovery, and unified multi-chain signing under a single framework wrapping both CGGMP21 and FROST. The integration with Ringtail post-quantum signatures provides a smooth upgrade path to quantum-resistant operations.

All source code is open-source at \texttt{github.com/luxfi/threshold/protocols/lss}.

%% ═══════════════════════════════════════════════════════════════════════════
\begin{thebibliography}{99}

\bibitem{shamir1979} A.~Shamir, ``How to Share a Secret,'' \emph{Communications of the ACM}, vol.~22, no.~11, pp.~612--613, 1979.

\bibitem{pedersen1991} T.~P.~Pedersen, ``Non-Interactive and Information-Theoretic Secure Verifiable Secret Sharing,'' in \emph{CRYPTO}, 1991.

\bibitem{feldman1987} P.~Feldman, ``A Practical Scheme for Non-Interactive Verifiable Secret Sharing,'' in \emph{FOCS}, 1987.

\bibitem{frost} C.~Komlo and I.~Goldberg, ``FROST: Flexible Round-Optimized Schnorr Threshold Signatures,'' in \emph{SAC}, 2020.

\bibitem{cggmp21} R.~Canetti \emph{et al.}, ``UC Non-Interactive, Proactive, Threshold ECDSA with Identifiable Aborts,'' in \emph{ACM CCS}, 2020.

\bibitem{ringtail} Z.~Kelling, ``Ringtail: Post-Quantum Lattice-Based Threshold Signatures,'' Lux Industries, Inc., 2024.

\bibitem{quasar} Z.~Kelling and V.~Seesahai, ``Quasar Consensus: Hybrid BLS + Ringtail Dual-Signature Finality,'' Lux Industries, Inc., 2022.

\bibitem{gjkr2003} R.~Gennaro, S.~Jarecki, H.~Krawczyk, and T.~Rabin, ``Secure Distributed Key Generation for Discrete-Log Based Cryptosystems,'' \emph{Journal of Cryptology}, vol.~20, no.~1, pp.~51--83, 2007.

\bibitem{bip340} P.~Wuille \emph{et al.}, ``BIP-340: Schnorr Signatures for secp256k1,'' 2020.

\bibitem{teleport} Z.~Kelling and V.~Seesahai, ``Lux Teleport: Sovereign Omnichain Liquidity,'' Lux Industries, Inc., 2022.

\end{thebibliography}

\end{document}
